\documentclass[11pt]{article}

\usepackage[margin=1in]{geometry}
\usepackage[T1]{fontenc}
\usepackage[utf8]{inputenc}
\usepackage{enumitem}
\usepackage{xurl}
\usepackage[hidelinks]{hyperref}

\setlist[enumerate]{leftmargin=*}

\begin{document}

\section{Scope}

The literature on EV adoption models is broad, but six model families are most
relevant for extending an affordance landscape model:

\begin{enumerate}
    \item Agent-based and microsimulation models.
    \item Diffusion and system dynamics models.
    \item Discrete choice and hybrid choice models.
    \item Social network and peer-effect models.
    \item Charging infrastructure and spatial adoption models.
    \item Policy and total cost of ownership models.
\end{enumerate}

\section{High-Level Findings}

EV adoption is usually not modeled as a single consumer-choice event. The
stronger models combine several mechanisms: household heterogeneity, price and
total-cost barriers, charging access, social visibility, policy incentives,
technology learning, and vehicle replacement timing. For an affordance model,
the most transferable idea is to treat EV adoption as the result of both
objective action opportunities and perceived feasibility.

The closest analogue to the current affordance model is the agent-based EV
adoption literature. Studies such as Eppstein et al. on PHEV market penetration,
McCoy and Lyons on Irish microsimulation, Brown's mixed-logit ABM, and Kangur et
al.'s psychologically grounded ABM all represent heterogeneous agents who update
or act under social and contextual influence. Useful sources include
Eppstein et al. \cite{eppstein2011}, McCoy and Lyons \cite{mccoy2014},
Brown \cite{brown2013}, and Kangur et al. \cite{kangur2017}.

Aggregate diffusion and system dynamics models are useful for feedback
structure, not micro-behavioral detail. They emphasize reinforcing loops among
adoption, infrastructure, cost reductions, legitimacy, and policy support.
Representative sources include
Struben and Sterman \cite{struben2008}, Shepherd et al. \cite{shepherd2012},
Pasaoglu et al. \cite{pasaoglu2016}, and the review by Gnann et al.
\cite{gnann2018}.

Discrete choice and hybrid choice models provide the strongest empirical
foundation for purchase probabilities. They estimate willingness to pay for
price, range, charging time, fuel/electricity cost, infrastructure, incentives,
and latent attitudes. Representative work includes Brownstone, Bunch, and Train
on joint RP/SP mixed logit \cite{brownstone2000}, Hidrue et al. on willingness
to pay \cite{hidrue2011}, Bolduc et al. on hybrid choice
\cite{bolduc2008}, Glerum et al. on hybrid choice \cite{glerum2014},
Helveston et al. on U.S. and China preferences \cite{helveston2015}, Kim et
al. on social influences and latent attitudes \cite{kim2014}, and Axsen et al.
on lifestyle heterogeneity \cite{axsen2015}.

Social network and peer-effect studies show that EV uptake is socially
interdependent. Peer adoption, local visibility, word of mouth, and perceived
normality shift the effective choice set. This maps directly onto social
learning in the affordance model. Useful sources include
McCoy and Lyons \cite{mccoy2014}, Kangur et al. \cite{kangur2017}, Liu et al.
\cite{liu2017}, Hsu et al. \cite{hsu2013}, and Edelenbosch et al.
\cite{edelenbosch2018}.

Charging infrastructure and spatial adoption models are central for an
affordance framing. Charging availability is not just an explanatory variable;
it is a local opportunity structure. Stronger studies distinguish charger
counts, accessibility, fast-charging corridors, home charging, land-use context,
and spatial spillovers. Representative sources include
Sierzchula et al. \cite{sierzchula2014}, Li et al. \cite{li2017},
Narassimhan and Johnson \cite{narassimhan2018}, White et al.
\cite{white2022}, and Qian et al. \cite{qian2024}.

Policy and total cost of ownership models are the best source for cost
mechanisms. They translate purchase subsidies, taxes, fuel prices, electricity
prices, charging costs, depreciation, maintenance, and vehicle lifetime into
adoption conditions. Useful sources include
Langbroek et al. \cite{langbroek2016}, Danielis et al.
\cite{danielis2018}, M\"{u}nzel et al. \cite{munzel2019}, and Noll et al.
\cite{noll2024}.

\section{Transferable Mechanisms For An Affordance-Based EV ABM}

The most useful model design is modular:

\begin{enumerate}
    \item Affordance layer: local charging access, home charging, public
    charging, fast-charging corridors, dealership/model availability, parking
    access, and financial affordability.
    \item Personal-state layer: environmental orientation, innovation affinity,
    risk tolerance, range anxiety, price sensitivity, and EV familiarity.
    \item Choice layer: probability of EV purchase conditional on replacement
    need, affordability, charging feasibility, and perceived usefulness.
    \item Social layer: peer exposure, neighborhood visibility, word of mouth,
    and social norm updating.
    \item Policy layer: purchase subsidies, taxes, fuel/electricity prices,
    charging infrastructure support, and information campaigns.
    \item Feedback layer: adoption increases visibility, encourages charger
    rollout, reduces perceived risk, and may reduce costs through technology
    learning.
\end{enumerate}

For the existing Mesa affordance model, the most natural extension is not to
replace the behavior rule with a full vehicle-choice model immediately. A safer
first step is to add EV-specific affordance fields and let them modify the
probability of pro-EV behavior:

\begin{verbatim}
effective_EV_propensity =
    personal_EV_state
    + peer_exposure_effect
    + charging_affordance_effect
    - price_or_TCO_barrier
    - range_anxiety_barrier
\end{verbatim}

This preserves the affordance logic while allowing empirical mechanisms from EV
adoption models to enter as interpretable modules.

\section{Specific Design Guidance}

For price mechanisms, use TCO literature to represent cost as more than upfront
price. Include purchase price, subsidy, expected fuel savings, electricity or
charging price, annual mileage, maintenance, home charger cost, and depreciation
where data are available. Price should affect feasibility and adoption
probability, not simply subtract from environmental attitude.

For infrastructure, avoid using only a global charger count. Treat charging as a
spatial affordance with home access, nearby public access, destination access,
fast-charging corridor access, reliability, and congestion if possible.

For social learning, use local adoption exposure or network-neighbor adoption to
change perceived feasibility, risk, and normality. Peer effects should be
allowed to amplify infrastructure and policy effects.

For heterogeneity, represent households by at least income or budget, dwelling
type, annual mileage, car replacement timing, environmental preference, and
innovation affinity. These are repeatedly important across choice, ABM, TCO, and
infrastructure studies.

\section{Key Gaps In The Literature}

Several gaps are relevant for a new research contribution:

\begin{enumerate}
    \item Many models still weakly connect micro-level social influence with
    spatial charging access.
    \item Used EV markets, resale value, vehicle availability, and supply
    constraints are under-modeled.
    \item Home charging, renter status, multifamily housing, and equity
    constraints are often too coarse.
    \item Public charging is often represented as station count, not
    reliability, congestion, price, or practical accessibility.
    \item Preference updating after trial, ownership experience, and peer
    information remains simplified.
    \item Policy interactions are usually clearer in system dynamics and TCO
    models than in household-level ABMs.
\end{enumerate}

\section*{References}

\begin{thebibliography}{99}

\bibitem{axsen2015}
Axsen, J., Bailey, J., \& Castro, M. A. (2015).
Preference and lifestyle heterogeneity among potential plug-in electric vehicle
buyers.
\emph{Energy Economics}, 50, 190--201.
\href{https://doi.org/10.1016/j.eneco.2015.05.003}{doi:10.1016/j.eneco.2015.05.003}.

\bibitem{bolduc2008}
Bolduc, D., Boucher, N., \& Alvarez-Daziano, R. (2008).
Hybrid choice modeling of new technologies for car choice in Canada.
\emph{Transportation Research Record}, 2082(1), 63--71.
\href{https://doi.org/10.3141/2082-08}{doi:10.3141/2082-08}.

\bibitem{brown2013}
Brown, M. (2013).
Catching the PHEVer: Simulating electric vehicle diffusion with an
agent-based mixed logit model of vehicle choice.
\emph{Journal of Artificial Societies and Social Simulation}, 16(2), 5.
\href{https://doi.org/10.18564/jasss.2127}{doi:10.18564/jasss.2127}.

\bibitem{brownstone2000}
Brownstone, D., Bunch, D. S., \& Train, K. (2000).
Joint mixed logit models of stated and revealed preferences for
alternative-fuel vehicles.
\emph{Transportation Research Part B: Methodological}, 34(5), 315--338.
\href{https://doi.org/10.1016/S0191-2615(99)00031-4}{doi:10.1016/S0191-2615(99)00031-4}.

\bibitem{danielis2018}
Danielis, R., Giansoldati, M., \& Rotaris, L. (2018).
A probabilistic total cost of ownership model to evaluate the current and future
prospects of electric cars uptake in Italy.
\emph{Energy Policy}, 119, 268--281.
\href{https://doi.org/10.1016/j.enpol.2018.04.024}{doi:10.1016/j.enpol.2018.04.024}.

\bibitem{edelenbosch2018}
Edelenbosch, O. Y., McCollum, D. L., Pettifor, H., Wilson, C., \& van Vuuren, D. P. (2018).
Interactions between social learning and technological learning in electric
vehicle futures.
\emph{Environmental Research Letters}, 13(12), 124004.
\href{https://doi.org/10.1088/1748-9326/aae948}{doi:10.1088/1748-9326/aae948}.

\bibitem{eppstein2011}
Eppstein, M. J., Grover, D. K., Marshall, J. S., \& Rizzo, D. M. (2011).
An agent-based model to study market penetration of plug-in hybrid electric
vehicles.
\emph{Energy Policy}, 39(6), 3789--3802.
\href{https://doi.org/10.1016/j.enpol.2011.04.007}{doi:10.1016/j.enpol.2011.04.007}.

\bibitem{glerum2014}
Glerum, A., Stankovikj, L., Th\'{e}mans, M., \& Bierlaire, M. (2014).
Forecasting the demand for electric vehicles: Accounting for attitudes and
perceptions.
\emph{Transportation Science}, 48(4), 483--499.
\href{https://doi.org/10.1287/trsc.2013.0487}{doi:10.1287/trsc.2013.0487}.

\bibitem{gnann2018}
Gnann, T., Stephens, T. S., Lin, Z., Pl\"{o}tz, P., Liu, C., \& Brokate, J. (2018).
What drives the market for plug-in electric vehicles? A review of international
PEV market diffusion models.
\emph{Renewable and Sustainable Energy Reviews}, 93, 158--164.
\href{https://doi.org/10.1016/j.rser.2018.03.055}{doi:10.1016/j.rser.2018.03.055}.

\bibitem{helveston2015}
Helveston, J. P., Liu, Y., Feit, E. M., Fuchs, E., Klampfl, E., \& Michalek, J. J. (2015).
Will subsidies drive electric vehicle adoption? Measuring consumer preferences
in the U.S. and China.
\emph{Transportation Research Part A: Policy and Practice}, 73, 96--112.
\href{https://doi.org/10.1016/j.tra.2015.01.002}{doi:10.1016/j.tra.2015.01.002}.

\bibitem{hidrue2011}
Hidrue, M. K., Parsons, G. R., Kempton, W., \& Gardner, M. P. (2011).
Willingness to pay for electric vehicles and their attributes.
\emph{Resource and Energy Economics}, 33(3), 686--705.
\href{https://doi.org/10.1016/j.reseneeco.2011.02.002}{doi:10.1016/j.reseneeco.2011.02.002}.

\bibitem{hsu2013}
Hsu, C.-I., Li, H.-C., \& Lu, S.-M. (2013).
A dynamic marketing model for hybrid electric vehicles: A case study of Taiwan.
\emph{Transportation Research Part D: Transport and Environment}, 20, 21--29.
\href{https://doi.org/10.1016/j.trd.2013.01.001}{doi:10.1016/j.trd.2013.01.001}.

\bibitem{kangur2017}
Kangur, A., Jager, W., Verbrugge, R., \& Bockarjova, M. (2017).
An agent-based model for diffusion of electric vehicles.
\emph{Journal of Environmental Psychology}, 52, 166--182.
\href{https://doi.org/10.1016/j.jenvp.2017.01.002}{doi:10.1016/j.jenvp.2017.01.002}.

\bibitem{kim2014}
Kim, J., Rasouli, S., \& Timmermans, H. (2014).
Expanding scope of hybrid choice models allowing for mixture of social
influences and latent attitudes: Application to intended purchase of electric
cars.
\emph{Transportation Research Part A: Policy and Practice}, 69, 71--85.
\href{https://doi.org/10.1016/j.tra.2014.08.016}{doi:10.1016/j.tra.2014.08.016}.

\bibitem{langbroek2016}
Langbroek, J. H. M., Franklin, J. P., \& Susilo, Y. O. (2016).
The effect of policy incentives on electric vehicle adoption.
\emph{Energy Policy}, 94, 94--103.
\href{https://doi.org/10.1016/j.enpol.2016.03.050}{doi:10.1016/j.enpol.2016.03.050}.

\bibitem{li2017}
Li, S., Tong, L., Xing, J., \& Zhou, Y. (2017).
The market for electric vehicles: Indirect network effects and policy design.
\emph{Journal of the Association of Environmental and Resource Economists},
4(1), 89--133.
\href{https://doi.org/10.1086/689702}{doi:10.1086/689702}.

\bibitem{liu2017}
Liu, C., Roberts, M. C., \& Sioshansi, R. (2017).
Spatial effects on hybrid electric vehicle adoption.
\emph{Transportation Research Part D: Transport and Environment}, 52, 85--97.
\href{https://doi.org/10.1016/j.trd.2017.02.014}{doi:10.1016/j.trd.2017.02.014}.

\bibitem{mccoy2014}
McCoy, D., \& Lyons, S. (2014).
Consumer preferences and the influence of networks in electric vehicle
diffusion: An agent-based microsimulation in Ireland.
\emph{Energy Research \& Social Science}, 3, 89--101.
\href{https://doi.org/10.1016/j.erss.2014.07.008}{doi:10.1016/j.erss.2014.07.008}.

\bibitem{munzel2019}
M\"{u}nzel, C., Pl\"{o}tz, P., Sprei, F., \& Gnann, T. (2019).
How large is the effect of financial incentives on electric vehicle sales? A
global review and European analysis.
\emph{Energy Economics}, 84, 104493.
\href{https://doi.org/10.1016/j.eneco.2019.104493}{doi:10.1016/j.eneco.2019.104493}.

\bibitem{narassimhan2018}
Narassimhan, E., \& Johnson, C. (2018).
The role of demand-side incentives and charging infrastructure on plug-in
electric vehicle adoption: Analysis of US states.
\emph{Environmental Research Letters}, 13(7), 074032.
\href{https://doi.org/10.1088/1748-9326/aad0f8}{doi:10.1088/1748-9326/aad0f8}.

\bibitem{noll2024}
Noll, B., Schmidt, T. S., \& Egli, F. (2024).
Managing trade-offs between electric vehicle taxation and adoption.
\emph{Cell Reports Sustainability}, 1(9), 100130.
\href{https://doi.org/10.1016/j.crsus.2024.100130}{doi:10.1016/j.crsus.2024.100130}.

\bibitem{pasaoglu2016}
Pasaoglu, G., Harrison, G., Jones, L., Hill, A., Beaudet, A., \& Thiel, C. (2016).
A system dynamics based market agent model simulating future powertrain
technology transition: Scenarios in the EU light duty vehicle road transport
sector.
\emph{Technological Forecasting and Social Change}, 104, 133--146.
\href{https://doi.org/10.1016/j.techfore.2015.11.028}{doi:10.1016/j.techfore.2015.11.028}.

\bibitem{qian2024}
Qian, L., Zhang, H., Chen, Y., \& Xiao, Y. (2024).
Spatial cross-side network effect of charging stations on electric vehicle
adoption.
\emph{Transportation Research Part D: Transport and Environment}, 133, 104400.
\href{https://doi.org/10.1016/j.trd.2024.104400}{doi:10.1016/j.trd.2024.104400}.

\bibitem{shepherd2012}
Shepherd, S., Bonsall, P., \& Harrison, G. (2012).
Factors affecting future demand for electric vehicles: A model based study.
\emph{Transport Policy}, 20, 62--74.
\href{https://doi.org/10.1016/j.tranpol.2011.12.006}{doi:10.1016/j.tranpol.2011.12.006}.

\bibitem{sierzchula2014}
Sierzchula, W., Bakker, S., Maat, K., \& van Wee, B. (2014).
The influence of financial incentives and other socio-economic factors on
electric vehicle adoption.
\emph{Energy Policy}, 68, 183--194.
\href{https://doi.org/10.1016/j.enpol.2014.01.043}{doi:10.1016/j.enpol.2014.01.043}.

\bibitem{struben2008}
Struben, J., \& Sterman, J. D. (2008).
Transition challenges for alternative fuel vehicle and transportation systems.
\emph{Environment and Planning B: Planning and Design}, 35(6), 1070--1097.
\href{https://doi.org/10.1068/b33022t}{doi:10.1068/b33022t}.

\bibitem{white2022}
White, L. V., Carrel, A. L., Shi, W., \& Sintov, N. D. (2022).
Why are charging stations associated with electric vehicle adoption? Untangling
effects in three United States metropolitan areas.
\emph{Energy Research \& Social Science}, 89, 102663.
\href{https://doi.org/10.1016/j.erss.2022.102663}{doi:10.1016/j.erss.2022.102663}.

\end{thebibliography}

\end{document}
